\section*{چکیده}

این گزارش درباره‌ی طراحی و بهینه‌سازی یک پایگاه‌داده \lr{PostgreSQL 16} برای مجموعه‌داده‌ی \lr{IMDb} است. این مجموعه‌داده خیلی بزرگ است -- حدود \lr{190} میلیون سطر در هفت جدول -- پس بارگذاری مستقیم آن هم کند است و هم فضای دیسک زیادی می‌گیرد. برای حل این مشکل، از یک \textbf{صف پیام} (\lr{message queue}) استفاده کردیم که با دستور \lr{SELECT ... FOR UPDATE SKIP LOCKED} کار می‌کند. این کار باعث می‌شود چند مصرف‌کننده هم‌زمان، بدون قفل‌شدن روی هم، داده را بنویسند. داده‌ها هم به‌صورت دسته‌ای (\lr{batch}) نوشته می‌شوند تا فشرده‌سازی \lr{TOAST} فعال شود و حجم روی دیسک کم شود.

بعد از بارگذاری، هشت کوئری تحلیلی که در سند پروژه خواسته شده بود را نوشتیم و اول بدون هیچ ایندکسی (خط پایه) با \lr{EXPLAIN (ANALYZE, BUFFERS)} بررسی کردیم. بعد شش ایندکس مناسب طراحی کردیم: چند \lr{B-Tree} جزئی برای فیلترهای بازه‌ای، و یک ایندکس \lr{GIN} ترکیبی (با کمک \lr{btree\_gin}) برای فیلتر روی \lr{title\_type}. طراحی ایندکس‌ها را با یک بازبینی خودکار چندعاملی بررسی کردیم و یک اسکریپت (\Rl{scripts/compare_timings.py}) هم نوشتیم که خودش جدول مقایسه‌ی «قبل/بعد» را از خروجی واقعی می‌سازد.

همه‌ی این پروژه -- بارگذاری، هشت سناریو، و ایندکس‌گذاری -- را واقعاً روی یک \lr{PostgreSQL 16} داخل \lr{Docker} اجرا کردیم. پس عددهای بخش \ref{sec:indexing} (زمان اجرا، تعداد بافر، نوع گره‌های پلن) واقعی‌اند، نه پیش‌بینی. فقط یک محدودیت وجود داشت: چون فضای دیسک سیستم اجرا کم بود، جدول \lr{title\_principals} فقط با \lr{8} میلیون از \lr{90} میلیون سطر واقعی‌اش بارگذاری شد (بخش \ref{sec:real-data-scale})؛ شش جدول دیگر کامل هستند. نتیجه هم یک‌دست نبود: شش سناریو سریع‌تر شدند، سناریوی \lr{3} تقریباً بدون تغییر ماند چون عمداً ایندکسی برایش نساختیم، و سناریوی \lr{1} با ایندکس \lr{GIN} حتی کندتر هم شد.
